%!TEX root = ../../thesis.tex

\section{Costs and Limitations}\label{sec:discussion-disadvantages}

% TODO: The counterweight to \cref{sec:discussion-advantages}, argued just as
% concretely and with numbers wherever they exist.
%
%   - Throughput: a conventional fuzzer executes orders of magnitude more inputs per
%     second. Give the actual gap from \cref{sec:platform-throughput} and state what
%     coverage the approach would have to reach per execution to break even.
%   - Reliability: collapse is a real and frequent outcome
%     (\cref{chap:failure-modes}), and a fuzzer that fails on a fraction of runs is a
%     hard sell against a technique that is boringly robust.
%   - Reward engineering: the strongest results are on targets whose reward was
%     shaped. That effort does not transfer, and it is exactly the manual per-target
%     work that made coverage-guided fuzzing win in the first place
%     (\cref{sec:fuzzing-limitations}).
%   - Action-space constraints: no masking, a flat categorical action and an
%     enumerated branching factor (\cref{sec:mdp-constraints}) rule out variable-length
%     inputs and phase-dependent legal actions without changing the learner.
%   - Observation scaling: a fixed hashed coverage vector aliases blocks and does not
%     grow with the target. A real stack has far more blocks than buckets.
%   - Engineering complexity: an emulator, a learner and a probe stack against a
%     fuzzer that is a single binary --- a cost worth stating plainly.
%
% Close by separating limitations of *this implementation* from limitations of the
% *approach*. Only the latter belong in \cref{chap:future-work} as open problems.
